\chapter[Local Comparison with Fisher]{Local Comparison with Fisher Information Geometry}
\section{Introduction}
The preceding chapters developed the intrinsic theory of the Universal
Optimality Defect (UOD). We now compare the Riemannian metric naturally
induced by the logarithmic embedding with the classical Fisher
information metric.
Let
$F=Q\circ L:\mathcal P^\circ\rightarrow F(\mathcal P^\circ),$
and
$\mathcal D(P,Q)=\|F(P)-F(Q)\|^2, \qquad g=F^{*}g_E.$
Throughout this chapter \(Q=I-\frac1n\mathbf{11}^T\) denotes the
orthogonal projection onto the zero-mean hyperplane.
\section{Local Quadratic Expansion}
\begin{theorem}[Local Quadratic Expansion]
\label{ch:thm:25-1}
For every \(P\in\mathcal P^\circ\) and every
\(\delta\in T_P\mathcal P^\circ\),
\[ \mathcal D(P+\delta,P) =
g_P(\delta,\delta) + o(\|\delta\|^2). \]
\end{theorem}
\begin{proof}
Taylor expansion of \(F\) gives
$F(P+\delta)=F(P)+DF_P(\delta)+o(\|\delta\|),$
therefore
\[ \mathcal D(P+\delta,P) =
\|DF_P(\delta)\|^2 + o(\|\delta\|^2). \]
Since \(g=F^{*}g_E\), the conclusion follows. 
\end{proof}
\section{Explicit Formula}
\begin{proposition}[Pullback Metric Formula]
\label{ch:prop:25-2}
For tangent vectors \(u,v\),
\[ g(u,v) = \sum_i\frac{u_iv_i}{P_i^2} -
\frac1n
\left(\sum_i\frac{u_i}{P_i}\right)
\left(\sum_i\frac{v_i}{P_i}\right). \]
\end{proposition}
\begin{proof}
Since
\[ DF_P=Q\circ DL_P, \qquad
DL_P(u)=\left(\frac{u_i}{P_i}\right), \]
and \(Q\) is orthogonal,
\[ g(u,v) = \langle Q\,DL_P u,Q\,DL_P v\rangle =
\langle DL_P u,Q\,DL_P v\rangle. \]
Expanding \(Q=I-\frac1n\mathbf{11}^T\) yields the formula. 
\end{proof}
\section{Fisher Metric}
The Fisher information metric is
$g_F(u,v) = \sum_i\frac{u_iv_i}{P_i}.$
\section{Matrix Representation}
Let
\[
D(P)=\operatorname{diag}\!\left(\frac1{P_i^2}\right),
\qquad v(P)=\left(\frac1{P_i}\right).
\]
Then
$G(P) = D(P)-\frac1n v(P)v(P)^T$
is the matrix representing the pullback metric.
The Fisher metric is represented by
\[
G_F(P)=\operatorname{diag}\!\left(\frac1{P_i}\right).
\]
\section{Uniform Equivalence}
\begin{theorem}[Uniform Metric Equivalence]
\label{ch:thm:25-3}
Let \(K\subset\mathcal P^\circ\) be compact.
Then there exist constants
$0<c_K\le C_K<\infty$
such that
$c_K\,g_F(u,u) \le g(u,u) \le C_K\,g_F(u,u)$
for every \(P\in K\) and every tangent vector
\(u\in T_P\mathcal P^\circ\).
\end{theorem}
\begin{proof}
Consider the normalized symmetric matrix
$A(P) = G_F(P)^{-1/2} G(P) G_F(P)^{-1/2}.$
Since both metrics are symmetric,
$g(u,u) = x^T A(P) x,$
where
$x = G_F(P)^{1/2}u.$
The matrix \(A(P)\) is symmetric and positive definite on the tangent
space. Hence, by the Rayleigh quotient theorem,
\[ \lambda_{\min}(A(P)) \|x\|^2 \le x^T A(P) x
\le \lambda_{\max}(A(P)) \|x\|^2. \]
Because
$\|x\|^2 = g_F(u,u),$
we obtain
\[ \lambda_{\min}(A(P))\,g_F(u,u) \le g(u,u)
\le \lambda_{\max}(A(P))\,g_F(u,u). \]
The entries of \(A(P)\) depend continuously on \(P\). Since \(K\) is compact,
its minimum and maximum eigenvalues are attained:
\[ c_K =
\min_{P\in K}\lambda_{\min}(A(P)),
\qquad C_K =
\max_{P\in K}\lambda_{\max}(A(P)). \]
Because \(A(P)\) is positive definite on the tangent space,
$0<c_K\le C_K<\infty.$
Therefore
$c_K\,g_F \le g \le C_K\,g_F.$
Consequently, the two metrics induce the same local topology on every
compact subset of the interior simplex. 
\end{proof}
\section{Consequences}
\begin{corollary}[Local Topological Equivalence]
\label{ch:cor:25-4}
On every compact subset:
\begin{itemize}
\item the pullback and Fisher metrics generate the same local topology;
\item they define equivalent notions of convergence;
\item infinitesimal estimates can be transferred between the two
    geometries up to uniform constants.
\end{itemize}
\end{corollary}
\section{Interpretation}
The Fisher metric arises from the Hessian of relative entropy, whereas
the WIS metric is obtained by transporting Euclidean geometry through
logarithmic quotient coordinates.
Although their global constructions differ, Theorem~\ref{ch:thm:25-3} shows that they
are locally equivalent on compact subsets of the posterior manifold.
\section{Outlook}
Having established the local relationship with Fisher geometry, we now
turn to global comparisons between the Universal Optimality Defect and
classical statistical divergences, beginning with Total Variation and
Hellinger distance.
